\begin{lexlean}{FoundryExperience}
\useglossary{lexlean.std.bool@1.2.0}
\useglossary{lexlean.std.nat@1.2.0}
\title{Boolean}
\begin{semanticmodule}
definition{
  name: oneLook,
  parameters: [],
  result: String,
  body: "Every surface Foundry draws itself, outside the renderer kit, takes the look of the rest of the app from the brand tokens in one stylesheet, in every theme: approval questions, notices of an open request, tool dialogs and the tools card. None writes a colour, a type face or a control style of its own. The result reads as clean, sharp and premium: one glass surface, a hairline edge, quiet text, no decoration."
}

definition{
  name: popupsCentred,
  parameters: [],
  result: String,
  body: "Every overlay popup is centred on the screen, across and down, over a light scrim. At most one approval question is shown at a time; the next waits its turn, and its time limit starts when it is shown."
}

definition{
  name: goldenRatio,
  parameters: [],
  result: String,
  body: "Every measure of these surfaces is from the golden ratio, 1.618: spaces from the series 5, 8, 13, 21, 34 and 55 pixels; a popup 550 pixels wide, which is 340 times the ratio, and at most 61.8 percent of the screen tall; 13 pixel text under a 16 pixel title; controls 34 pixels tall, and at least 44 under a finger."
}

definition{
  name: questionAsNotification,
  parameters: [],
  result: String,
  body: "A build-time question is drawn as the app draws its own notifications: a mark in a soft tile, a title, the explanation, the exact details behind one row, and two answers at the right. The refusal is plain and holds the focus; the approval is solid and is the last control."
}

definition{
  name: whyHowWhat,
  parameters: [],
  result: String,
  body: "Every explanation is ordered why, then how, then what: why this is being asked, how it happens and where the data goes, then what exactly is being approved and what cannot be undone. It is concise and self-descriptive, in plain words without technical jargon, and it still states every fact the approval boundaries require to be disclosed."
}

definition{
  name: sidebarIsNavigation,
  parameters: [],
  result: String,
  body: "The sidebar holds only navigation and is 233 pixels wide, which is 144 times the ratio."
}

definition{
  name: toolsOnSettings,
  parameters: [],
  result: String,
  body: "The tools are the first card of the Settings page: a title, then one quiet control each, with its own mark."
}

definition{
  name: homeAsksOneThing,
  parameters: [],
  result: String,
  body: "Home asks one thing, what to build: it shows the question alone, with no line under it, no app picker and no mode selector. Every chat builds; the mode selector is not shown, focused or offered in Settings."
}

definition{
  name: homeIsAWorkspace,
  parameters: [],
  result: String,
  body: "Home is a workspace, not a page. Its line, Build your ideas with Prism, is one row at every size, the sentence in the page's own face and the name in a serif italic, so the line reads as a sentence and the name as a signature. The hero fills the first screen at every size, and on a desktop the floor of apps already made peeks 89 pixels above the bottom edge, so one short scroll brings the whole row and its way into the gallery into view. On a phone home is one screen that never scrolls: the floor gives way to the Apps tab, which already lists every app. The hero sits at the golden section of the room above the floor, its question over a room at most 987 pixels wide; the box it is asked in is that room's golden section, 610 pixels, centred, with three chips centred under it on one line: two ideas, a fresh two on every visit, and Surprise me, which writes another into the box. On a tall window the room under the hero is capped at 233 pixels, so it settles toward the middle instead of leaving a void. The floor is one row on a panel that rises from the bottom edge, inset from the sides, on the popover glass under a hairline, with a card 233 by 144 for each app and its name under it. The floor is looked at four ways, on one track at its left: Search, Featured, My apps and Bookmarked, the chosen one lit, with the way to all of them at its right. Featured is the same four Prism apps for every visitor, whether or not they have made anything; My apps is what this browser holds; Bookmarked is what the visitor starred, kept in this browser alone and never sent anywhere. A card is its cover, a word on the cover saying whether it is published, then its name and who made it under its name. The cards share the floor's width in even columns, four on a wide screen and two on a narrow one, so the row fills the floor at any size. Each app's cover is drawn on the card, white light entering and leaving as a spectrum turned to that app's own hue, so nothing is fetched to show it."
}

definition{
  name: composerRow,
  parameters: [],
  result: String,
  body: "The box's row reads left to right as the question is made: attach, then the model that will answer beside it, then the three ways to ask, then send. The way of asking is chosen on the same pill as the model, with its name and a chevron, and its menu lists each way on two lines, its name over a few words on what choosing it means. The names are the ones the kernel's catalogue gives: Build, Chat and Plan."
}

definition{
  name: featuredCovers,
  parameters: [],
  result: String,
  body: "Each app on the floor shows its own picture, taken from the app itself the first time it runs and kept until the app changes, filling the card from its top edge. An app that has not run yet shows a tile in the brand's light carrying its initial, so every card reads as a card and none is an empty box."
}

definition{
  name: modelMenu,
  parameters: [],
  result: String,
  body: "The model menu lists each model on two lines, its name over a few words on what choosing it means, with a tag saying where it runs, Cloud or Local, each with its mark, in one column at the right."
}

definition{
  name: appsOnDevice,
  parameters: [],
  result: String,
  body: "The record of what apps do on this device leads with one sentence that sums every app up: how many opened, how many reached the internet and how many were stopped. Each app is then one line, its name and what it reached, that opens to each site it was allowed or stopped from reaching, how often and when, the space it takes and when it last opened; a dot says the same at a glance. It says app, site, opened, reached, allowed and stopped, never the network terms behind them. The explanation answers why, then how, then what without naming them, and how fast builds are is one line that opens on demand."
}
\end{semanticmodule}
\end{lexlean}
